



\section{Algorithm for a variant of the Uhlmann transformation in the mixed sample access model}
\label{sec:altanative mixed sample uhlmann}


 

We provide a quantum algorithm for implementing a variant of the Uhlmann transformation, given mixed sample access to $\rho^\sA$ and $\sigma^\sA$.
The variant transformation corresponds to a partial isometry $V^\sA$ such that
\begin{align}
\label{inteq:144}
    \rF(\rho^\sA, \sigma^\sA) = \rF(V^\sA\ket{\rho_{\rm c}}^{\sA\sB}, \ket{\sigma_{\rm c}}^{\sA\sB}),
\end{align}
where $\ket{\rho_{\rm c}}^{\sA\sB}$ and $\ket{\sigma_{\rm c}}^{\sA\sB}$ are the canonical purified states of $\rho^\sA$ and $\sigma^\sA$, respectively.
Recall that this partial isometry $V^\sA$ is equivalent to the Uhlmann partial isometry constructed from states $\rho^\top$ and $\sigma^\top$ on the system $\sB$.
As shown by the inequality:
\begin{align}
    \rF(\rho^\sA, \sigma^\sA) 
    &= \rF(V^\sA\ket{\rho_{\rm c}}^{\sA\sB}, \ket{\sigma_{\rm c}}^{\sA\sB}) \\
    &\leq \rF\big(V^\sA\rho^\sA(V^\sA)^\dag, \sigma^\sA\big),
\end{align}
the partial isometry $V^\sA$ has the property that it always brings $\rho^\sA$ closer to $\sigma^\sA$.


Our statement is as follows, where we denote by $\rho_{\rm min}$ and $\sigma_{\rm min}$ the minimum non-zero eigenvalues of $\rho^\sA$ and $\sigma^{\sA}$, respectively, and by $s_{\rm min}$ and $r$ the minimum non-zero singular value and the rank of $\sqrt{\sigma^\sA}\sqrt{\rho^\sA}$, respectively.
It should be noted that $d_\sA = d_\sB$. 





\begin{theorem}[Algorithm for a variant of the Uhlmann transformation in the mixed sample access model]
\label{thm:variant of uhlmann mixed sample}
    
Let $\delta \in (0, 1)$ and $\chi \in \{\diamond, \tF\}$. 
Then, the quantum sample algorithm $\mathtt{VariantUhlmannMixedSample}$ given by \Cref{alg:Uthm:variant of uhlmann mixed sample} satisfies the following.

A quantum channel $\cJ_\diamond^{\sA\sH}$ given by $\cJ_\diamond^{\sA\sH}  
    =\mathtt{VariantUhlmannMixedSample}(\rho^{\hat{\sA}}, \sigma^{\hat{\sA}}; \delta, \diamond)$, satisfies that
\begin{equation}
    \f{1}{2}\big\|\cJ_\diamond^{\sA\sH}\circ\cP_{\ket{0}}^{\bC\rarr\sH} - \cU_{\rm ideal}^{\sA\sH}\circ\cP_{\ket{0}}^{\bC\rarr\sH}\big\|_\diamond \leq \delta,
\end{equation}
where $\til{U}_{\rm ideal}^{\sA\sH}$ is an exact block-encoding unitary of $V^\sA$, and $V^\sA$ is a partial isometry such that
\begin{align}
    \rF(V^\sA\ket{\rho_{\rm c}}^{\sA\sB}, \ket{\sigma_{\rm c}}^{\sA\sB}) = \rF(\rho^\sA, \sigma^\sA).
\end{align}
The algorithm uses $\zeta_\diamond$ samples of $\rho^{\hat{\sA}}$ and $\sigma^{\hat{\sA}}$, where 
\begin{align}
\label{inteq:146}
    \zeta_\diamond = \til{\cO}\Big(\f{1}{\delta s_{\rm min}^2} \min\Big\{\f{1}{\rho_{\rm min}^2} + \f{1}{\sigma_{\rm min}^2}, \f{1}{\delta^4 s_{\rm min}^4}\Big\}\Big).
\end{align}

A quantum channel $\cJ_\tF^{\sA\sH}$ given by $\cJ_\tF^{\sA\sH} 
    = \mathtt{VariantUhlmannMixedSample}(\rho^{\hat{\sA}}, \sigma^{\hat{\sA}}; \delta, \tF)$, satisfies that
\begin{align}
    \rF\big(\cT^{\sA}(\ketbra{\rho_{\rm c}}{\rho_{\rm c}}^{\sA\sB}), \ket{\sigma_{\rm c}}^{\sA\sB}\big) \geq \rF(\rho^\sA, \sigma^\sA) - \delta,
\end{align}
where $\cT^{\sA}= \tr_\sH\circ\cJ_\tF^{\sA\sH}\circ\cP_{\ket{0}}^{\bC\rarr\sH}$.
The algorithm uses $\zeta_{\tF}$ samples of $\rho^{\hat{\sA}}$ and $\sigma^{\hat{\sA}}$, where 
\begin{align}
\label{inteq:145}
    \zeta_{\tF} = \til{\cO}\Big(\f{1}{\delta\beta_\tF}\min\Big\{\f{1}{\rho_{\rm min}^2} + \f{1}{\sigma_{\rm min}^2}, \f{1}{\beta_\tF^4\delta^4}\Big\}\Big),
\end{align}
and $\beta_\tF = \f{1}{8}\max\{s_{\rm min}, \delta/(2r)\}$.

In both cases, the quantum circuit for implementing the algorithm consists of $\cO\big(\zeta_\chi \log{d_\sA}\big)$ one- and two-qubit gates, and $\cO\big(\log{d_\sA}\big)$ qubits suffice at any one time.


\end{theorem}




Unlike \Cref{alg:Uhl mixed sample} in Sec.~\ref{sec:mixed sample}, \Cref{alg:Uthm:variant of uhlmann mixed sample} can be implemented with substantially fewer samples, since Eqs.~\eqref{inteq:146} and~\eqref{inteq:145} do not include the dimensional factor $d_\sA$, which appears in the sample complexity in Theorem~\ref{thm:Uhlmann mixed state sample}.
This is because \Cref{alg:Uthm:variant of uhlmann mixed sample} avoids the use of the canonical purification algorithm $\mathtt{CanonicalPurification}$ in \Cref{alg:canonical purification}, which takes a high sample complexity.
Instead, it employs $\mathtt{BrockEncSqrtState}$ in \Cref{prop:block enc sqrt state} to directly block-encode the square root of given mixed states.





\begin{algorithm}[h]
\caption{Algorithm for a variant of the Uhlmann transformation in the mixed sample access model \parbox{\linewidth}{\centering $\mathtt{VariantUhlmannMixedSample}(\rho^{\hat{\sA}}, \sigma^{\hat{\sA}}; \delta, \chi)$ (In Theorem~\ref{thm:variant of uhlmann mixed sample})}}
\label{alg:Uthm:variant of uhlmann mixed sample}
\SetKwInput{KwInput}{Input}
\SetKwInput{KwOutput}{Output}
\SetKwInput{KwParameters}{Parameters}
\SetKwComment{Comment}{$\triangleright$\ }{}
\SetCommentSty{textnormal}
\SetKwProg{Fn}{Subroutine}{}{end}
\SetKwFunction{UPS}{UhlmannPurifiedSample}

\SetAlgoNoEnd
\SetAlgoNoLine
\KwInput{Two quantum states $\rho^{\hat{\sA}}$ and $\sigma^{\hat{\sA}}$.}
\KwParameters{$\delta \in (0, 1)$ and $\chi \in \{\diamond, \tF\}$.}
\KwOutput{Quantum channel $\cJ^{\sA\sH}$.}
\SetAlgoLined

\If{$\chi = \diamond$}{
  Set $\delta_2 \gets (\delta/6)^2$ and $\beta \gets s_{\rm min}/8$. \\
}
\ElseIf{$\chi = \tF$}{
  Set $\delta_2 \gets \delta/8$ and $\beta \gets \max\{s_{\rm min}/8, \delta_2/(2r)\}$. \\
}
Set $u$ to be the minimum odd integer satisfying $u \geq \big\lceil\f{8e}{\beta}\log{(2/\delta_2)}\big\rceil$. \\
Set $\delta_1 \gets \delta/(4u)$. \\
Set $\cF^{\hat{\sA}\sC_1} \gets \mathtt{BlockEncSqrtState}(\rho^{\hat{\sA}}; \delta_1)$ 
and $\cF_{\rm inv}^{\hat{\sA}\sC_1} \gets \mathtt{BlockEncSqrtState}^\dag(\rho^{\hat{\sA}}; \delta_1)$. \\
Set $\cG^{\hat{\sA}\sC_2} \gets \mathtt{BlockEncSqrtState}(\sigma^{\hat{\sA}}; \delta_1)$ 
and $\cG_{\rm inv}^{\hat{\sA}\sC_2} \gets \mathtt{BlockEncSqrtState}^\dag(\sigma^{\hat{\sA}}; \delta_1)$ (Proposition~\ref{prop:block enc sqrt state}). \\
Set $\cL^{\hat{\sA}\sC_1\sC_2} \gets \cG^{\hat{\sA}\sC_2} \circ \cF^{\hat{\sA}\sC_1}$
and $\cL_{\rm inv}^{\hat{\sA}\sC_1\sC_2} \gets \cF_{\rm inv}^{\hat{\sA}\sC_1} \circ \cG_{\rm inv}^{\hat{\sA}\sC_2}$.
Set $\cJ^{\sA\sH} \gets \mathtt{QSVTSIGN}(\cL^{\hat{\sA}\sC_1\sC_2}, \cL_{\rm inv}^{\hat{\sA}\sC_1\sC_2}; \beta, \delta_2)$ (Proposition~\ref{prop:FPAA general}). \\
Return $\cJ^{\sA\sH}$.
\end{algorithm}





\begin{proof}[Proof of Theorem~\ref{thm:variant of uhlmann mixed sample}]
    
From the Uhlmann's theorem,  there exists a partial isometry $V^\sA$ which satisfies
\begin{align}
\label{inteq:49}
    \rF(V^\sA\ket{\rho_{\rm c}}^{\sA\sB}, \ket{\sigma_{\rm c}}^{\sA\sB}) 
    &= \rF(\rho_{\rm c}^{\sB}, \sigma_{\rm c}^{\sB}) \\
    &= \rF\big((\rho^\sA)^\top, (\sigma^\sA)^\top\big) \\
    \label{inteq:129}
    &= \rF(\rho^\sA, \sigma^\sA).
\end{align}
An explicit form of this partial isometry $V^\sA$ is given by 
\begin{align}
    V^\sA 
    &= \sgn^{(\rm SV)}\big(\tr_\sB\big[\ketbra{\sigma_{\rm c}}{\rho_{\rm c}}^{\sA\sB}\big]\big) \\
    &= \sgn^{(\rm SV)}\big(\tr_\sB\big[\sqrt{\sigma^\sA}\ketbra{\Gamma}{\Gamma}^{\sA\sB}\sqrt{\rho^\sA}\big]\big) \\
    &= \sgn^{(\rm SV)}\big(\sqrt{\sigma^\sA}\sqrt{\rho^\sA}\big),
\end{align}
where $\ket{\Gamma}^{\sA\sB} = \sum_i\ket{i}^\sA\ket{i}^\sB$.
Thus, it suffices to directly construct a unitary that block-encodes $\sqrt{\sigma^\sA}\sqrt{\rho^\sA}$, and then lift up all singular values to unity using the QSVT for sign function $\mathtt{QSVTSIGN}$ (Proposition~\ref{prop:FPAA general}).


From a block-encoding of the square root of a quantum state (Proposition~\ref{prop:block enc sqrt state}), we obtain quantum channels $\cF^{\hat{\sA}\sC_1} = \mathtt{BlockEncSqrtState}(\rho^{\hat{\sA}}; \delta_1)$ and $\cG^{\hat{\sA}\sC_2} = \mathtt{BlockEncSqrtState}(\sigma^{\hat{\sA}}; \delta_1)$, such that
\begin{align}
\label{inteq:127}
    \f{1}{2}\big\|\cF^{\hat{\sA}\sC_1} - \cW_{\sqrt{\rho}}^{\hat{\sA}\sC_1}\big\|_\diamond \leq \delta_1, \ \ \ \text{and} \ \ \
    \f{1}{2}\big\|\cG^{\hat{\sA}\sC_2} - \cW_{\sqrt{\sigma}}^{\hat{\sA}\sC_2}\big\|_\diamond \leq \delta_1,
\end{align}
where $W_{\sqrt{\rho}}^{\hat{\sA}\sC_1}$ and $W_{\sqrt{\sigma}}^{\hat{\sA}\sC_2}$ are $(1, 5, 0)$-block-encoding unitary of $\sqrt{\rho^{\hat{\sA}}}/(2\sqrt{2})$ and $\sqrt{\sigma^{\hat{\sA}}}/(2\sqrt{2})$, respectively, using $h = \til{\cO}\Big(\f{1}{\delta_1}\min\Big\{\f{1}{\rho_{\rm min}^2} + \f{1}{\sigma_{\rm min}^2}, \f{1}{\delta_1^4}\Big\}\Big)$, samples of $\rho^{\hat{\sA}}$ and $\sigma^{\hat{\sA}}$.

 
We here use the following lemma regarding the product of block-encoded matrices.
\begin{lemma}[Product of block-encoded matrices~{\cite[Lemma 30]{gilyen2019qsvt}}; see also Ref.~{\cite[Lemma 3.3.10]{Gilyn2019thesis}}]
\label{lem:product of BE}
Suppose that $U_1^{\sA\sB}$ is an $(\alpha_1, a_1, \epsilon_1)$-block-encoding unitary of a matrix $M_1^\sA$, and $U_2^{\sA\sC}$ is an $(\alpha_2, a_2, \epsilon_2)$-block-encoding unitary of a matrix $M_2^\sA$. Then, $(\bI^\sC\otimes U_1^{\sA\sB})(\bI^\sB \otimes U_2^{\sA\sC})$ is an $(\alpha_1\alpha_2, a_1+a_2, \alpha_1\epsilon_2 + \alpha_2\epsilon_1)$-block-encoding of $M_1^\sA M_2^\sA$.
\end{lemma}

Let $W^{\sA\sC_1\sC_2} = W_{\sqrt{\sigma}}^{\hat{\sA}\sC_2}W_{\sqrt{\rho}}^{\hat{\sA}\sC_1}$.
From \Cref{lem:product of BE}, we see that
\begin{align}
    W^{\sA\sC_1\sC_2} 
    = 
\begin{blockarray}{ccc}
 & \bra{0}^{\sC_1\sC_2} &  & \vspace{1mm}\\
\begin{block}{c(cc)}
  \ket{0}^{\sC_1\sC_2} \hspace*{1mm} & \f{1}{8}\sqrt{\sigma^\sA}\sqrt{\rho^\sA} & \hspace{0mm} * \hspace{3mm}\\
   \hspace*{1mm} & * & \hspace{0mm} * \hspace{3mm}\\
\end{block}
\end{blockarray}\hspace{2.5mm},
\end{align}
and when we set $\cL^{\hat{\sA}\sC_1\sC_2} = \cG^{\hat{\sA}\sC_2}\circ\cF^{\hat{\sA}\sC_1}$, we have 
\begin{align}
\label{inteq:130}
    \f{1}{2}\big\|\cL^{\hat{\sA}\sC_1\sC_2} - \cW^{\hat{\sA}\sC_1\sC_2}\big\|_\diamond \leq 2\delta_1,
\end{align}
where we used Eq.~\eqref{inteq:127}.
Thus, $(1,10,0)$-block-encoding unitary $W^{\hat{\sA}\sC_1\sC_2}$ of $\sqrt{\sigma^{\hat{\sA}}}\sqrt{\rho^{\hat{\sA}}}/8$ is approximately obtained with error $2\delta_1$.
Similarly, a quantum channel $\cL_{\rm inv}^{\hat{\sA}\sC_1\sC_2}$ that satisfies $\f{1}{2}\|\cL_{\rm inv}^{\hat{\sA}\sC_1\sC_2} - (\cW^{\hat{\sA}\sC_1\sC_2})^\dag\|_\diamond \leq 2\delta_1$ is obtained from the same number of samples of $\rho^{\hat{\sA}}$ and $\sigma^{\hat{\sA}}$.


The rest is similar to the discussion in Sec.~\ref{sec:proofUhlmann}.
Let $\cJ^{\sA\sH} = \mathtt{QSVTSIGN}(\cL^{\hat{\sA}\sC_1\sC_2}, \cL_{\rm inv}^{\hat{\sA}\sC_1\sC_2}; \beta, \delta_2)$ and $u = \cO\big(\log(1/\delta_2)/\beta\big)$.
First, regarding the evaluation in the diamond norm, we set $\beta = \beta_\diamond = s_{\rm min}/8$, and then, it holds that $\big\|P_\sgn^{(\rm SV)}(\til{M}^\sA) - V^\sA\big\|_\infty \leq \delta_2$, where $\til{M}^\sA = \sqrt{\sigma^{\hat{\sA}}}\sqrt{\rho^{\hat{\sA}}}/8$ and $V^\sA = \sgn^{(\rm SV)}(\til{M}^\sA)$ that correspond to the partial isometry.
From \Cref{lem:FPAA diamond}, we have
\begin{align}
\label{inteq:128}
    \f{1}{2}\big\|\cJ^{\sA\sH}\circ\cP_{\ket{0}}^{\bC\rarr\sH} - \cU_{\rm ideal}^{\sA\sH} \circ \cP_{\ket{0}}^{\bC\rarr\sH}\big\|_\diamond \leq 3\sqrt{\delta_2} + 2u\delta_1.
\end{align}
Rescaling the parameters as $\delta_1 = \delta/(4u)$ and $\delta_2 = (\delta/6)^2$, Eq.~\eqref{inteq:128} is bounded by $\delta$. In this case, the number of samples of $\rho^{\hat{\sA}}$ and $\sigma^{\hat{\sA}}$ are evaluated as $\zeta_\diamond = hu$, that is, 
\begin{align}
    \zeta_\diamond 
    &= \cO\Big(\f{1}{\beta_\diamond}\log{\Big(\f{1}{\delta_2}\Big)}\Big)\til{\cO}\Big(\f{1}{\delta_1}\min\Big\{\f{1}{\rho_{\rm min}^2} + \f{1}{\sigma_{\rm min}^2}, \f{1}{\delta_1^4}\Big\}\Big) \\
    &= \til{\cO}\Big(\f{1}{\delta s_{\rm min}^2}\min\Big\{\f{1}{\rho^2}+\f{1}{\sigma_{\rm min}^2}, \f{1}{\delta^4s_{\rm min}^4}\Big\}\Big).
\end{align}



Next, we evaluate the error in fidelity difference.
Let $\cT^\sA = \tr_\sH\circ \cJ^{\sA\sH} \circ \cP_{\ket{0}}^{\bC\rarr\sH}$.
From an almost identical calculation to that in Sec.~\ref{sec:uhlfidpurifsamp}, by setting $\beta = \beta_\tF = \max\{s_{\rm min}/8, \delta_2/(2r)\}$, we have that
\begin{align}
    \rF\big(\cT^\sA(\ketbra{\rho_{\rm c}}{\rho_{\rm c}}^{\sA\sB}), \ket{\sigma_{\rm c}}^{\sA\sB}\big)
    &\geq \rF\big(V^\sA\ket{\rho_{\rm c}}^{\sA\sB}, \ket{\sigma_{\rm c}}^{\sA\sB}\big) - \f{u}{2}\big\|\cL^{\hat{\sA}\sC_1\sC_2} - \cW^{\hat{\sA}\sC_1\sC_2}\big\|_\diamond -4\delta_2 \\
    &\geq \rF(\rho^\sA, \sigma^\sA) - 2u\delta_1 - 4\delta_2,
\end{align}
where we used Eqs.~\eqref{inteq:129} and~\eqref{inteq:130}.
Hence, rescaling the parameters as $\delta_1 = \delta/(4u)$ and $\delta_2 = \delta/8$ yields
the result that $\rF\big(\cT^\sA(\ketbra{\rho_{\rm c}}{\rho_{\rm c}}^{\sA\sB}), \ket{\sigma_{\rm c}}^{\sA\sB}\big) \geq \rF(\rho^\sA, \sigma^\sB) -\delta$.
The number of samples $\zeta_\tF = hu$ of $\rho^{\hat{\sA}}$ and $\sigma^{\hat{\sA}}$ is
\begin{align}
    \zeta_\tF 
    &= \cO\Big(\f{1}{\beta_\tF}\log{\Big(\f{1}{\delta_2}\Big)}\Big)\til{\cO}\Big(\f{1}{\delta_1}\min\Big\{\f{1}{\rho_{\rm min}^2} + \f{1}{\sigma_{\rm min}^2}, \f{1}{\delta_1^4}\Big\}\Big) \\
    &=\til{\cO}\Big(\f{1}{\delta\beta_\tF^2}\min\Big\{\f{1}{\rho_{\rm min}^2} + \f{1}{\sigma_{\rm min}^2}, \f{1}{\delta^4\beta_\tF^4}\Big\}\Big),
\end{align}
where $\beta_\tF = \f{1}{8}\max\{s_{\rm min}, \delta/(2r)\}$.

The algorithm $\mathtt{BlockEncSqrtState}$ uses $\cO(h\log{d_\sA})$ one- and two-qubit gates and is repeated $\cO(u)$ times in $\mathtt{QSVTSIGN}$.
Thus, the total number of gates in the quantum circuit for the entire algorithm is $\cO(hu\log{d_\sA}) = \cO(\zeta_\chi\log{d_\sA})$, where $\chi \in \{\diamond, \tF\}$.
Due to the sequential property of the algorithm, $\cO(\log{d_\sA})$ qubits suffice at any one time.



\end{proof}















%=========================================================================================================================================================



\section{Naive state tomography-based approach for the Uhlmann transformation}
\label{sec:Comparing with a naive tomography-based}



The most straightforward strategy for implementing the Uhlmann transformation is to use quantum state tomography~\cite{ODonnell2016EffTomo, Haah2017samploptTomo, Guta2020FastTomo, Apeldoorn2023TomoPurifQuery, Chen2023WhenAdapTomo, hu2024sampleoptimalmemoryefficient}.
We first obtain approximate classical descriptions of $\ket{\rho}^{\sA\sB}$ and $\ket{\sigma}^{\sA\sB}$ via state tomography, and then directly compute the Uhlmann partial isometry $V^\sB = \sgn^{(\rm SV)}\big(\tr_{\sA}\big[\ketbra{\sigma}{\rho}^{\sA\sB}\big]\big)$ on a classical computer.
This partial isometry coincides with that determined by the Schmidt bases of $\ket{\rho}^{\sA\sB}$ and $\ket{\sigma}^{\sA\sB}$.
To implement the partial isometry by quantum circuits, we need to extend it to a unitary and, for instance, decompose it into quantum circuits by a brute-force method. This clearly results in an exponential number of gates in general~\cite{nielsen2010quantum}.
Moreover, in a classical computation step, this approach generally requires exponential time and space in terms of classical bits.



Below, we provide a statement about the query and sample complexities of the Uhlmann transformation based on the state tomography approach.
Here, the results in the purified and mixed sample access models can be obtained by allowing collective measurements over multiple copies, potentially even an exponential number of them.
The derivations are discussed in Appendices~\ref{sec:tomography based purif query},~\ref{sec:tomography based purif sample}, and~\ref{sec:tomography based mixed sample}, corresponding to each of the three computational models.


\begin{proposition}[Uhlman transformation by a state tomography-based approach]
    In each of the three query/sample models, for $\delta \in (0, 1)$, there exists a quantum query/sample algorithm based on state tomography that outputs a partial isometry $\til{V}^{\sB}$ satisfying $\big\|\til{V}^\sB - V^\sB\big\|_\infty \leq \delta$, where $V^\sB$ is the Uhlmann partial isometry.
    The query/sample complexity of the algorithm in each model is given as follows:
\begin{itemize}
\item Purified query access model: $\til{\cO}\Big(\f{d_\sA d_\sB}{\delta s_{\rm min}}\Big)$ queries to $U_\rho^{\sA\sB}$, $U_\sigma^{\sA\sB}$, and their inverses.
\item Purified sample access model: $\til{\cO}\Big(\f{d_\sA d_\sB}{\delta^2 s_{\rm min}^2}\Big)$ samples of $\ket{\rho}^{\sA\sB}$ and $\ket{\sigma}^{\sA\sB}$.
\item Mixed sample access model: $\til{\cO}\Big(\f{d_\sA (r_\rho + r_\sigma)}{\delta^4 s_{\rm min}^4}\Big)$ samples of $\rho^\sA$ and $\sigma^\sA$.
\end{itemize}
The algorithms generally require exponential time and space in classical computation, as well as an exponential number of one- and two-qubit gates.
\end{proposition}







\subsection{In the purified query access model}
\label{sec:tomography based purif query}





In the purified query access model, the best-known query complexity for state tomography of a pure state in a $d$-dimensional Hilbert space, up to trace distance error $\epsilon$, is given by $\til{\cO}(d/\epsilon)$~\cite{Apeldoorn2023TomoPurifQuery}.
Thus, to obtain a classical description of $\til{\rho}^{\sA\sB}$ and $\til{\sigma}^{\sA\sB}$ that satisfy
\begin{align}
\label{eq:tomography descp. error}
    \f{1}{2}\|\til{\rho}^{\sA\sB} - \ketbra{\rho}{\rho}^{\sA\sB}\|_1 \leq \epsilon, 
    \ \ \text{and} \ \ \  
    \f{1}{2}\|\til{\sigma}^{\sA\sB} - \ketbra{\sigma}{\sigma}^{\sA\sB}\|_1 \leq \epsilon,
\end{align}
it suffice to make $\til{\cO}\big(d_\sA d_\sB/\epsilon\big)$ queries to the unitaries $U_\rho^{\sA\sB}$ and $U_\sigma^{\sA\sB}$, which prepare $\ket{\rho}^{\sA\sB}$ and $\ket{\sigma}^{\sA\sB}$, respectively.
We then compute an approximation $\til{V}^\sB$ of the Uhlmann partial isometry $V^\sB$ from the classical description of $\til{\rho}^{\sA\sB}$ and $\til{\sigma}^{\sA\sB}$.


While there is some choice in how to construct $\til{V}^\sB$ from $\til{\rho}^{\sA\sB}$ and $\til{\sigma}^{\sA\sB}$, we specifically consider $\til{V}^\sB = \sgn^{(\rm SV)}\big(\tr_\sA[\ketbra{\til{\sigma}_1}{\til{\rho}_1}^{\sA\sB}]\big)$, where $\ket{\til{\rho}_1}^{\sA\sB}$ and $\ket{\til{\sigma}_1}^{\sA\sB}$ are eigenstates corresponding to the largest eigenvalues of $\til{\rho}^{\sA\sB}$ and $\til{\sigma}^{\sA\sB}$, respectively.
As we will see below, this construction yields a pretty good approximation of $V^\sB$, supported by results from matrix perturbation theory in mathematical physics.


To evaluate how well $\til{V}^\sB$ approximates the desired partial isometry $V^\sB$ up to a global phase, we analyze the spectral properties. Let the eigenvalue decomposition of $\til{\rho}^{\sA\sB}$ be given by $\til{\rho}^{\sA\sB} = \sum_{j=1}^{r_{\til{\rho}}} \til{\rho}_j \ketbra{\til{\rho}_j}{\til{\rho}_j}^{\sA\sB}$, where $\til{\rho}_1 \geq \til{\rho}_2 \geq \ldots \geq \til{\rho}_{r_{\til{\rho}}}$.
First, to evaluate the distance between $\ket{\til{\rho}_1}^{\sA\sB}$ and $\ket{\rho}^{\sA\sB}$, we use a special case of the Weyl-type matrix perturbation theorem~\cite{Weyl1912DasAV, mirsky1960symmetric, Li1998perturvationEigSing, bhatia2013matrix}.


\begin{lemma}[Matrix spectrum perturbation on the Schatten-$p$ norm~{\cite[Theorem 5]{mirsky1960symmetric}}]
\label{lem:weyl tipe perturb}
Let $A$ and $\til{A}$ be $n_1 \times n_2$ matrices with singular values $s_1(A) \geq \ldots \geq s_{\min\{n_1, n_2\}}(A)$ and $s_1(\til{A}) \geq \ldots \geq s_{\min\{n_1, n_2\}}(\til{A})$, respectively. Then, for any $p \in [1, \infty]$, we have $\big\|{\rm diag}\big(s_j(A) - s_j(\til{A})\big)\big\|_p \leq \big\|A-\til{A}\big\|_p$ where ${\rm diag}(a_j)$ is a diagonal matrix whose diagonal elements are $a_j$.
\end{lemma}

From this lemma of $p=1$ and Eq.~\eqref{eq:tomography descp. error}, we see that
\begin{align}
\label{inteq:85}
    |1 - \til{\rho}_1| + \sum_{j=2}^{r_{\til{\rho}}} |\rho_j| \leq 2\epsilon.
\end{align}
Thus, we have 
\begin{align}
    \big\|\ketbra{\til{\rho}_1}{\til{\rho}_1}^{\sA\sB} - \ketbra{\rho}{\rho}^{\sA\sB}\big\|_1 
    &\leq \big\|\ketbra{\til{\rho}_1}{\til{\rho}_1}^{\sA\sB} - \til{\rho}^{\sA\sB}\big\|_1 + \big\|\til{\rho}^{\sA\sB} - \ketbra{\rho}{\rho}^{\sA\sB}\big\|_1 \\
    &= |1 - \til{\rho}_1| + \sum_{j=2}^{r_{\til{\rho}}} |\rho_j| + \big\|\til{\rho}^{\sA\sB} - \ketbra{\rho}{\rho}^{\sA\sB}\big\|_1 \\
    \label{inteq:86}
    &\leq 4 \epsilon,
\end{align}
where we used Eqs.~\eqref{eq:tomography descp. error} and~\eqref{inteq:85}.
Hence, for the state $\til{\rho}^{\sA\sB}$ obtained by state tomography, its eigenvector $\ket{\til{\rho}_1}^{\sA\sB}$ corresponding to the largest eigenvalue sufficiently approximates $\ket{\rho}^{\sA\sB}$.
The same holds for the states $\ket{\til{\sigma}_1}^{\sA\sB}$ and $\ket{\sigma}^{\sA\sB}$:
\begin{align}
    \label{inteq:107}
    \big\|\ketbra{\til{\sigma}_1}{\til{\sigma}_1}^{\sA\sB} - \ketbra{\sigma}{\sigma}^{\sA\sB}\big\|_1
    &\leq 4 \epsilon.
\end{align}

Next, we evaluate the distance between the matrices before applying the sign function, i.e., $\tr_\sA\big[\ketbra{\til{\sigma}_1}{\til{\rho}_1}^{\sA\sB}\big]$ and $\tr_\sA\big[\ketbra{\sigma}{\rho}^{\sA\sB}\big]$.
It is important to note that there is a degree of freedom associated with the global phase $e^{i\theta}$ on the matrix, which does not affect the final result. 
We here fix the phase as $\theta = -\theta_\rho + \theta_\sigma$, where $e^{-i\theta_\rho} = \braket{\til{\rho}_1}{\rho} / |\braket{\til{\rho}_1}{\rho}|$ and $e^{-i\theta_\sigma} = \braket{\til{\sigma}_1}{\sigma}/|\braket{\til{\sigma}_1}{\sigma}|$.
These phases achieve the minimum in the relation between the trace norm and the Euclidean norm (see also Eq.~\eqref{eq:relation of Euclidean and trace norm}).

Let $\til{M}^\sB = \tr_{\sA}\big[\ketbra{\til{\sigma}_1}{\til{\rho}_1}^{\sA\sB}\big]$ and $M^\sB = e^{i\theta}\tr_{\sA}\big[\ketbra{\sigma}{\rho}^{\sA\sB}\big]$.
We see that
\begin{align}
    \big\|\til{M}^\sB - M^\sB\big\|_\infty
    &\leq \big\|\til{M}^\sB - M^\sB\big\|_1 \\
    &= \big\|\tr_{\sA}\big[\ketbra{\til{\sigma}_1}{\til{\rho}_1}^{\sA\sB}\big] - e^{i\theta}\tr_{\sA}\big[\ketbra{\sigma}{\rho}^{\sA\sB}\big]\big\|_1 \\
    &\leq \big\|\ketbra{\til{\sigma}_1}{\til{\rho}_1}^{\sA\sB}
    - e^{i\theta}\ketbra{\sigma}{\rho}^{\sA\sB}\big\|_1 \\
    &\leq \big\|\ket{\til{\sigma}_1}^{\sA\sB}(\bra{\til{\rho}_1}^{\sA\sB} - e^{-i\theta_\rho}\bra{\rho}^{\sA\sB}) \big\|_1 +  \big\|(\ket{\til{\sigma}_1}^{\sA\sB} - e^{i\theta_\sigma}\ket{\sigma}^{\sA\sB})
    \bra{\rho}^{\sA\sB}\big\|_1 \\
    &\leq \big\|\ket{\til{\rho}_1}^{\sA\sB} - e^{i\theta_\rho}\ket{\rho}^{\sA\sB}\big\|
    + \big\|\ket{\til{\sigma}_1}^{\sA\sB} - e^{i\theta_\sigma}\ket{\sigma}^{\sA\sB}\big\| \\
    &\leq \f{1}{\sqrt{2}}\big\|\ketbra{\til{\rho}_1}{\til{\rho}_1}^{\sA\sB} - \ketbra{\rho}{\rho}^{\sA\sB}\big\|_1 + \f{1}{\sqrt{2}}\big\|\ketbra{\til{\sigma}_1}{\til{\sigma}_1}^{\sA\sB} - \ketbra{\sigma}{\sigma}^{\sA\sB}\big\|_1    \\
    \label{inteq:37}
    &\leq 4\sqrt{2}\epsilon,
\end{align}
where we used Eq.~\eqref{eq:relation of Euclidean and trace norm} in the fifth inequality, and used Eqs.~\eqref{inteq:86} and~\eqref{inteq:107} in the last inequality.
Thus, we can well approximate the matrix $M^\sB$ within the error $4\sqrt{2}\epsilon$.




Finally, to evaluate the distance between $\til{V}^\sB$ and $V^\sB$, we again refer to a matrix perturbation technique. 
By the polar decomposition, any matrices $A$ and $\til{A}$ can be decomposed as $A = U\sqrt{AA^\dag}$ and $\til{A} = \til{U}\sqrt{\til{A}\til{A}^\dag}$, where $U$ and $\til{U}$ are partial isometry polar factors of $A$ and $\til{A}$, respectively. 
The distance between the partial isometry polar factors has been studied in Refs.~\cite{Li1993AperturbGen, Li2002PertBound, Li2005SomeNewPerturb, XiaoShan2008VariationsQHfactor, Liu2008PerturbApprox, Li2008SubunitaryPolar, Zhang2013NewPerturb, Hong2014SomePertuPolarDecomp, Duong2016EffectOfPerturb, zhu2018NewPerturbationUniInv, Fu2020AnOptimallPerturb}.
To the best of our knowledge, the tightest bound for the operator norm of these partial isometries without any rank constraint is given in Ref.~\cite{Zhang2013NewPerturb}:
\begin{align}
\label{inteq:39}
    \big\|\til{U} - U\big\|_\infty 
    \leq \sqrt{\Big(\f{2}{a_{\rm min} + \til{a}_{\rm min}}\Big)^2 + \f{1}{a_{\rm min}^2} + \f{1}{\til{a}_{\rm min}^2}} \big\|A-\til{A}\big\|_\infty,
\end{align}
where $a_{\rm min}$ and $\til{a}_{\rm min}$ are the minimum non-zero singular values of $A$ and $\til{A}$, respectively.

We apply this to evaluating $\big\|\til{V}^\sB -V^\sB\big\|_\infty$.
The matrix $M^\sB$ and $\til{M}^\sB$ are decomposed as $M^\sB = (V^\sB)^\dag\sqrt{M^\sB(M^\sB)^\dag}$ and $\til{M}^\sB = (\til{V}^\sB)^\dag\sqrt{\til{M}^\sB(\til{M}^\sB)^\dag}$, since $V^\sB=\sgn^{(\rm SV)}(M^\sB)$ and $\til{V}^\sB=\sgn^{(\rm SV)}(\til{M}^\sB)$.
We then obtain 
\begin{align}
\label{inteq:38}
    \big\|\til{V}^\sB - V^\sB\big\|_\infty
    &= \big\|(\til{V}^\sB)^\dag - (V^\sB)^\dag\big\|_\infty \\
    &\leq \f{\sqrt{3}\big\|\til{M}^\sB - M^\sB\big\|_\infty}{\min\big\{s_{\rm min}(M), s_{\rm min}(\til{M})\big\}} \\
    \label{inteq:87}
    &\leq \f{4\sqrt{6}\epsilon}{s_{\rm min}(M) - 4\sqrt{2}\epsilon},
\end{align}
where we used Eqs.~\eqref{inteq:37} and~\eqref{inteq:39}, and used $|s_{\rm min}(M) - s_{\rm min}(\til{M})| \leq 4\sqrt{2}\epsilon$ which is a result from \Cref{lem:weyl tipe perturb} of $p=\infty$ and Eq.~\eqref{inteq:37}.
Thus, in order to approximate $V^\sB$ within error $\delta$, it is sufficient to choose the error $\epsilon$ in state tomography such that $\epsilon = \f{\delta s_{\rm min}(M)}{4\sqrt{2}(\sqrt{3} + \delta)}$.
Note that with this choice, $s_{\rm min}(M) - 4\sqrt{2}\epsilon > 0$ holds.




Consequently, we conclude that to satisfy $\big\|\til{V}^\sB - V^\sB\big\|_\infty \leq \delta$, the required query complexity is given by $\til{\cO}\Big(\f{d_\sA d_\sB}{\epsilon}\Big) 
    = \til{\cO}\Big(\f{d_\sA d_\sB}{\delta s_{\rm min}}\Big)$, where we used the fact that $s_{\rm min}(M)$ is the same as the minimum non-zero singular value of $\sqrt{\sigma^\sA}\sqrt{\rho^\sA}$, i.e., $s_{\rm min}$.
This completes the derivation of the query complexity for the Uhlmann transformation using the naive state tomography-based approach.


While we derived only an upper bound, there is little hope for a substantial improvement in this approach, because both the state tomography and the bound in Eq.~\eqref{inteq:39} are known to be nearly optimal.
Moreover, the minimum non-zero singular values that appear in Eq.~\eqref{inteq:39} are not the result of $U$ and $\til{U}$ being partial isometries. A similar bound is obtained even in the case that $U$ and $\til{U}$ are full-rank, i.e., unitaries~\cite{chun1989perturbation, Barrlund1990PerturbOnPolar, Mathias1993perturbPDecomp, Li1995NewPerturbUnitary, bhatia2013matrix}. 






%================================================================================================================================






\subsection{In the purified sample model}
\label{sec:tomography based purif sample}

The quantum state tomography-based approach in the purified sample access model is mostly the same as the approach in the purified query access model described in \Cref{sec:tomography based purif query}.
The difference from the discussion in \Cref{sec:tomography based purif query} lies only in the number of samples required in state tomography.
The optimal sample complexity for a pure state in a $d$-dimensional Hilbert space, up to trace distance error $\epsilon$, is given by $\til{\cO}(d/\epsilon^2)$~\cite{ODonnell2016EffTomo, Haah2017samploptTomo}, where collective measurements over multiple copies, potentially even exponentially many, are allowed.
Thus, to obtain a classical description of $\til{\rho}^{\sA\sB}$ and $\til{\sigma}^{\sA\sB}$ that satisfy 
\begin{align}
\label{eq:tomography descp. error 2}
    \f{1}{2}\|\til{\rho}^{\sA\sB} - \ketbra{\rho}{\rho}^{\sA\sB}\|_1 \leq \epsilon, \ \ \text{and} \ \ \ 
    \f{1}{2}\|\til{\sigma}^{\sA\sB} - \ketbra{\sigma}{\sigma}^{\sA\sB}\|_1 \leq \epsilon,
\end{align}
we use $\til{\cO}\big(d_\sA d_\sB/\epsilon^2\big)$ samples of $\ket{\rho}^{\sA\sB}$ and~$\ket{\sigma}^{\sA\sB}$.

Following the same strategy in \Cref{sec:tomography based purif query}, it follows that for implementing the Uhlmann transformation within the error $\delta$, the total number of samples of $\ket{\rho}^{\sA\sB}$ and $\ket{\sigma}^{\sA\sB}$ is $\til{\cO}\Big(\f{d_\sA d_\sB}{\epsilon^2}\Big) 
    = \til{\cO}\Big(\f{d_\sA d_\sB}{\delta^2 s_{\rm min}^2}\Big)$.






%==================================================================================================================================



\subsection{In the mixed sample access model}
\label{sec:tomography based mixed sample}


In the mixed sample access model, from multiple copies of $\rho^\sA$ and $\sigma^\sA$, we obtain approximations $\til{\rho}^\sA$ and $\til{\sigma}^\sA$ of the original states, respectively.
The optimal sample complexity of state tomography for a rank-$k$ quantum state in a $d$-dimensional Hilbert space, up to a trace distance error of $\epsilon$, is given by $\til{\cO}(kd/\epsilon^2)$~\cite{ODonnell2016EffTomo, Haah2017samploptTomo},
where one is allowed to perform collective measurements on multiple copies, possibly even an exponential number of them.




Unlike the other two models, the mixed sample access model requires more careful analysis.
Specifically, even though $\omega^\sA$ can be estimated with a small error, it is necessary to evaluate whether the purified state can also be estimated with a small error.
In the evaluation, we use the following inequality: for any state $\til{\omega}^\sA$ and $\omega^\sA$,
\begin{align}
\label{inteq:26}
    \f{1}{2}\big\|\ketbra{\til{\omega}_{\rm c}}{\til{\omega}_{\rm c}}^{\sA\sB} - \ketbra{\omega_{\rm c}}{\omega_{\rm c}}^{\sA\sB}\big\|_1 
    \leq \big\|\til{\omega}^\sA - \omega^\sA\big\|_1^{1/2},
\end{align}
where $\ket{\omega_{\rm c}}^{\sA\sB}$ and $\ket{\til{\omega}_{\rm c}}^{\sA\sB}$ are the canonical purified states of $\omega^\sA$ and $\til{\omega}^\sA$, respectively, and $d_\sA = d_\sB$.
We derive this inequality at the end of this section.

Since the optimal sample complexity to obtain $\til{\rho}^\sA$ and $\til{\sigma}^\sA$ which satisfy
\begin{align}
    \f{1}{2}\big\|\til{\rho}^\sA - \rho^\sA \big\|_1 \leq \epsilon',
    \ \  \text{and} \ \ \  
    \f{1}{2}\big\|\til{\sigma}^\sA - \sigma^\sA \big\|_1 \leq \epsilon',
\end{align}
are given by $\til{\cO}\big(d_\sA r_\rho/\epsilon'^2\big)$ and $\til{\cO}\big(d_\sA r_\sigma /\epsilon'^2\big)$, respectively, where $r_\rho$ and $r_\sigma$ are the ranks of $\rho$ and $\sigma$.

Once classical descriptions of $\til{\rho}^\sA$ and $\til{\sigma}^\sA$ are obtained, one can simply diagonalize them via classical computation, which immediately yields classical descriptions of $\ketbra{\til{\rho}_{\rm c}}{\til{\rho}_{\rm c}}^{\sA\sB}$ and $\ketbra{\til{\sigma}_{\rm c}}{\til{\sigma}_{\rm c}}^{\sA\sB}$.
By Eq.~\eqref{inteq:26}, to obtain classical descriptions of $\ketbra{\til{\rho}_{\rm c}}{\til{\rho}_{\rm c}}^{\sA\sB}$ and $\ketbra{\til{\sigma}_{\rm c}}{\til{\sigma}_{\rm c}}^{\sA\sB}$ such that 
\begin{align}
    \f{1}{2}\big\|\ketbra{\til{\rho}_{\rm c}}{\til{\rho}_{\rm c}}^{\sA\sB} - \ketbra{\rho_{\rm c}}{\rho_{\rm c}}^{\sA\sB}\big\|_1 \leq \epsilon, \ \  \text{and} \ \ \ 
    \f{1}{2}\big\|\ketbra{\til{\sigma}_{\rm c}}{\til{\sigma}_{\rm c}}^{\sA\sB} - \ketbra{\sigma_{\rm c}}{\sigma_{\rm c}}^{\sA\sB}\big\|_1 \leq \epsilon,
\end{align}
the number of samples of $\rho^\sA$ and $\sigma^\sA$ is $\til{\cO}(d_\sA (r_\rho + r_\sigma) / \epsilon^4)$, where $\epsilon'$ is chosen as $\epsilon' = \epsilon^2$.



The procedure hereafter follows similarly to that in \Cref{sec:tomography based purif query}.
We classically compute $\til{V}^{\sB} = \sgn^{(\rm SV)}\big(\tr_\sA\big[\ketbra{\til{\sigma}_{\rm c}}{\til{\rho}_{\rm c}}^{\sA\sB}\big]\big)$ from the description of $\ket{\til{\rho}_{\rm c}}^{\sA\sB}$ and $\ket{\til{\sigma}_{\rm c}}^{\sA\sB}$. This partial isometry $\til{V}^{\sB}$ can be a good approximation of the ideal Uhlmann partial isometry $V^{\sB} = \sgn^{(\rm SV)}\big(\tr_\sA\big[\ketbra{\sigma_{\rm c}}{\rho_{\rm c}}^{\sA\sB}\big]\big)$.
The total number of samples of $\rho^\sA$ and $\sigma^\sA$ for the Uhlmann transformation within the error $\delta$ using the state tomography-based approach, is given by 
\begin{align}
    \til{\cO}\Big(\f{d_\sA (r_\rho + r_\sigma)}{\epsilon^4}\Big) 
    \label{eq:tomobase sample mixed}
    = \til{\cO}\Big(\f{d_\sA (r_\rho + r_\sigma)}{\delta^4 s_{\rm min}^4}\Big).
\end{align}



As an alternative approach, one might compute $\big((\sqrt{\til{\rho}}\sqrt{\til{\sigma}})^{\sB}\big)^\top$ directly, since it holds that
\begin{align}
    \tr_\sA\big[\ketbra{\sigma_{\rm c}}{\rho_{\rm c}}^{\sA\sB}\big] 
    &= \tr_\sA\big[\sqrt{\sigma^\sA}\ketbra{\Gamma}{\Gamma}^{\sA\sB}\sqrt{\rho^\sA}\big] \\
    &= \Big(\big(\sqrt{\rho}\sqrt{\sigma}\big)^{\sB}\Big)^\top.
\end{align}
In this case, we see that
\begin{align}
    \big\|\big(\sqrt{\til{\rho}}\sqrt{\til{\sigma}}\big)^\top - \big(\sqrt{\rho}\sqrt{\sigma}\big)^\top\big\|_\infty
    &=\big\|\sqrt{\til{\rho}}\sqrt{\til{\sigma}} - \sqrt{\rho}\sqrt{\sigma}\big\|_\infty \\
    &\leq\|\sqrt{\til{\rho}} - \sqrt{\rho}\|_\infty 
    + \|\sqrt{\til{\sigma}} - \sqrt{\sigma}\|_\infty \\
    &\leq\|\sqrt{\til{\rho}} - \sqrt{\rho}\|_2 
    + \|\sqrt{\til{\sigma}} - \sqrt{\sigma}\|_2 \\
    &\leq\|\til{\rho} - \rho\|_1^{1/2} 
    + \|\til{\sigma} - \sigma\|_1^{1/2} \\
    \label{inteq:88}
    &\leq 2\sqrt{2\epsilon'},
\end{align}
where, in the third inequality we used the Powers--St\o rmer inequality in Eq.~\eqref{eq:powers stormer ineq}.
We find that Eq.~\eqref{inteq:88} has the same error scaling as in the case where we compute $\tr_\sA\big[\ketbra{\sigma_{\rm c}}{\rho_{\rm c}}^{\sA\sB}\big]$. 
Since the tomography error $\epsilon'$ is chosen as $\epsilon' = \epsilon^2$ and the singular values of $(\sqrt{\rho}\sqrt{\sigma})^\top$ coincide with those of $\sqrt{\sigma}\sqrt{\rho}$, the sample complexity is equivalent to that given in Eq.~\eqref{eq:tomobase sample mixed}.







Finally, we derived the inequality in Eq.~\eqref{inteq:26} via a technique of vectorization, which is represented by a linear map $\mathrm{Vec}$ such that, for a given orthonormal basis $\{\ket{i}\}_i$, $\mathrm{Vec}(\ket{\psi}\bra{\varphi}) = \ket{\psi}\ket{\varphi^*}$, where the complex conjugate is taken in the basis $\{\ket{i}\}_i$, i.e., $\ket{\varphi^*} = \sum_ic_i^*\ket{i}$ when $\ket{\varphi} = \sum_i c_i\ket{i}$.
It is straightforward to verify that the vectorization has the property that, for any matrix $M$, $\|M\|_2 = \|\mathrm{Vec}(M)\|$.

We utilize this property in our evaluation. This implies that
\begin{align}
    \big\|\sqrt{\til{\omega}^\sA} - \sqrt{\omega^\sA}\big\|_2 
    &= \big\|\mathrm{Vec}\big(\sqrt{\til{\omega}^\sA}\big) - \mathrm{Vec}\big(\sqrt{\omega^\sA}\big)\big\|\\
    &= \big\|\ket{\til{\omega}_{\rm c}}^{\sA\sB} - \ket{\omega_{\rm c}}^{\sA\sB}\big\| \\
    \label{inteq:27}
    &\geq \f{1}{2}\big\|\ketbra{\til{\omega}_{\rm c}}{\til{\omega}_{\rm c}}^{\sA\sB} - \ketbra{\omega_{\rm c}}{\omega_{\rm c}}^{\sA\sB}\big\|_1,
\end{align}
where we used that $\ket{\omega_{\rm c}}^{\sA\sB} = \mathrm{Vec}\big(\sqrt{\omega^\sA}\big)$ in the second equation, and used the relation between the trace norm and the Euclidean norm in Eq.~\eqref{eq:relation of Euclidean and trace norm} in the inequality.
Moreover, using the Powers--St\o rmer inequality in Eq.~\eqref{eq:powers stormer ineq}, we obtain that
\begin{align}
\label{inteq:28}
    \big\|\sqrt{\til{\omega}^\sA} - \sqrt{\omega^\sA}\big\|_2 \leq \big\|\til{\omega}^\sA - \omega^\sA\big\|_1^{1/2}.
\end{align}
From Eqs.~\eqref{inteq:27} and~\eqref{inteq:28}, we complete the derivation of Eq.~\eqref{inteq:26}.







%=============================================================================================================================================================================









\section{Overview of previous approach to the Uhlmann transformation}
\label{sec:metger and yuen's algorithm}

An explicit algorithm for the Uhlmann transformation was proposed by T. Metger and H. Yuen in Ref.~\cite{metger2023stateqipspace}.
This algorithm demonstrates that the Uhlmann transformation can be implemented with polynomial space complexity.
We here provide a brief and high-level overview of the algorithm tailored to the purified query access model; for full details, see Ref.~\cite{metger2023stateqipspace}.





\begin{theorem}[Space efficient Uhlmann transformation algorithm in the purified query access model~{\cite[Theorem 7.4 in its full version]{metger2023stateqipspace}}]
\label{prevthm:Metger&Yuen's result}
For $\delta \in (0, 1)$, there exists a quantum query algorithm that realizes a quantum channel $\cT^\sB$ satisfying 
\begin{align}
    \rF\big(\cT^\sB(\ketbra{\rho}{\rho}^{\sA\sB}), \ket{\sigma}^{\sA\sB}\big) \geq \rF(\rho^\sA, \sigma^\sA) - \delta,
\end{align}
using $\gamma$ queries to $U_\rho^{\sA\sB}$ and $U_\sigma^{\sA\sB}$, and their inverses, where $\gamma = \til{\cO}\Big(\f{d_\sA^5 d_\sB^5}{\delta^2}\min\Big\{\f{d_\sA^9 d_\sB^6}{s_{\rm min}^3}, \f{r^3}{\delta^3}\Big\}\Big)$. The quantum circuit of this algorithm consists of $\til{\cO}(\gamma)$ one- and two-qubit gates, and $\cO\big(\log{(d_\sA d_\sB)}\big)$ qubits suffice at any one time.
\end{theorem}



We should note that, at every stage of this algorithm, even in classical computation, only polynomial space is used; it is designed so that an exponential amount of classical data is not stored.
Meanwhile, as seen from the subsequent derivation, this approach involves operations such as incrementally summing measurement outcomes, which generally leads to exponential time in the number of bits in classical computation.





In the rest of this appendix, we derive \Cref{prevthm:Metger&Yuen's result}.
A main idea behind this algorithm is sketched as follows.
In the first step, we prepare the state $\ket{\rho}^{\sA\sB} = U_\rho^{\sA\sB}\ket{0}^{\sA\sB}$.
Then, we apply Pauli measurements to estimate $\tr[\Lambda^{\sA\sB}\rho^{\sA\sB}]$ for all Pauli operators $\Lambda^{\sA\sB} \in \{\bI, X, Y, Z\}^{\otimes d_{\sA\sB}^2}$, where $\rho^{\sA\sB} = \ketbra{\rho}{\rho}^{\sA\sB}$.
From the estimated value $\til{\alpha}_\Lambda$, we classically compute an element $\til{c}_{ij} = \f{1}{d_{\sA\sB}}\sum_\Lambda \til{\alpha}_\Lambda \bra{i}\Lambda\ket{j}^{\sA\sB}$, and subsequently prepare a two-qubit gate
\begin{align}
\til{U}_{ij}^{\sC} = 
\begin{pmatrix} 
\til{c}_{ij} & * \ \ \ \\
\sqrt{1 - |\til{c}_{ij}|^2} & * \ \ \ 
\end{pmatrix}.
\end{align}
Then, we construct a unitary $W_\rho^{\sA\sB\hat{\sA}\hat{\sB}}$ such that
\begin{align}
    W_\rho^{\sA\sB\hat{\sA}\hat{\sB}\sC} = H^{\sA\sB} \big( \sum_{i, j} \ketbra{j}{j}^{\sA\sB} \otimes \ketbra{i}{i}^{\hat{\sA}\hat{\sB}} \otimes \til{U}_{ij}^\sC \big) H^{\hat{\sA}\hat{\sB}},
\end{align}
where $H$ is the Hadamard operator: $H^\sA\ket{i}^\sA = \f{1}{\sqrt{d_\sA}}\sum_j (-1)^{i\cdot j}\ket{j}^\sA$.
We can see that this unitary satisfies $\bra{0}^{\hat{\sA}\hat{\sB}\sC}W_\rho^{\sA\sB\hat{\sA}\hat{\sB}\sC}\ket{0}^{\hat{\sA}\hat{\sB}\sC} = \til{\rho}^{\sA\sB}/d_{\sA\sB}$, where $\til{\rho}^{\sA\sB} = \f{1}{d_{\sA\sB}}\sum_\Lambda \til{\alpha}_\Lambda \Lambda^{\sA\sB}$.

Let $\sD = \hat{\sA}\hat{\sB}\sC$.
We analyze the number of queries required for $W_\rho^{\sA\sB\sD}$ to be a good approximation of a block-encoding of $\rho^{\sA\sB}$.
For all $d_{\sA\sB}^2$ Pauli operators $\Lambda^{\sA\sB}$, to obtain $\til{\alpha}_\Lambda$ such that $\big|\til{\alpha}_\Lambda - \tr[\Lambda^{\sA\sB}\rho^{\sA\sB}]\big| \leq \delta_1/d_{\sA\sB}$ with probability at least $1 - \eta$, it suffices to have $d_{\sA\sB}^2\cO\big(\f{d_{\sA\sB}^2}{\delta_1^2} \log(d_{\sA\sB}^2/\eta)\big) = \til{\cO}(d_{\sA\sB}^4/\delta_1^2)$ copies of $\ket{\rho}^{\sA\sB} = U_\rho^{\sA\sB}\ket{0}^{\sA\sB}$~\cite{aaronson2007learnability, yuen2022lecturetomography}. This is the result from the Chernoff bound for i.i.d. sampling~\cite{chernoff1952measure, hoeffding1963Probability, motwani1995randomized, mitzenmacher2005probability}.
By sequentially summing over each $\Lambda^{\sA\sB}$, the value $\til{c}_{ij}$ can be computed using only polynomial space.
Since this procedure is performed for all $i$ and $j$, the total number of queris is given by $d_{\sA\sB}^2 \til{\cO}(d_{\sA\sB}^4/\delta_1^2) = \til{\cO}(d_{\sA\sB}^6/\delta_1^2)$.
At this time, we have
\begin{align}
    \big\|d_{\sA\sB}\bra{0}^{\sD}W_\rho^{\sA\sB\sD}\ket{0}^{\sD} - \rho^{\sA\sB}\big\|_\infty 
    &=\big\|\til{\rho}^{\sA\sB} - \rho^{\sA\sB}\big\|_\infty \\
    &\leq \f{1}{d_{\sA\sB}}\sum_\Lambda |\til{\alpha}_\Lambda - \tr[\Lambda^{\sA\sB}\rho^{\sA\sB}]\big| \big\|\Lambda^{\sA\sB}\big\|_\infty \\
    &\leq \f{1}{d_{\sA\sB}}\sum_\Lambda |\til{\alpha}_\Lambda - \tr[\Lambda^{\sA\sB}\rho^{\sA\sB}]\big| \\
    &\leq \delta_1.
\end{align}
Hence, $W_\rho^{\sA\sB\sD}$ is an $(d_{\sA\sB}, 1 + \log{d_{\sA\sB}}, \delta_1)$-block-encoding of $\rho^{\sA\sB}$.
Similarly, we perform this procedure for $\ket{\sigma}^{\sA\sB} = U_\sigma^{\sA\sB}\ket{0}^{\sA\sB}$ and obtain an approximate block-encoding unitary $W_\sigma^{\sA\sD}$ from the same number of queris.


In the next step, we construct a unitary 
\begin{align}
    U^{\sA'\sB\sD_1\sD_2} = W_\sigma^{\sA'\sB\sD_1}(U_\sigma^{\sA'\sB})^\dag U_\rho^{\sA'\sB}(W_\rho^{\sA'\sB\sD_2})^\dag,
\end{align}
which is a $(d_{\sA\sB}^2, 2(1+\log{d_{\sA\sB}}), 2d_{\sA\sB}\delta_1)$-block-encoding unitary of $\ketbra{\sigma}{\rho}^{\sA'\sB}$.
To show this, we use \Cref{lem:product of BE} regarding the product of block-encoded matrices.
The unitaries $W_\rho^{\sA'\sB\sD_2}(U_\rho^{\sA'\sB})^\dag$ and $W_\sigma^{\sA'\sB\sD_1}(U_\sigma^{\sA'\sB})^\dag$ are $(d_{\sA\sB}, 1 + \log{d_{\sA\sB}}, \delta_1)$-block-encoding unitaries of $\ketbra{\rho}{0}^{\sA'\sB}$ and $\ketbra{\sigma}{0}^{\sA'\sB}$, respectively.
The claim that $U^{\sA'\sB\sD_1\sD_2}$ is $(d_{\sA\sB}^2, 2(1+\log{d_{\sA\sB}}), 2d_{\sA\sB}\delta_1)$-block-encoding unitary of $\ketbra{\sigma}{\rho}^{\sA'\sB}$ follows from \Cref{lem:product of BE}.

We need to trace over system $\sA'$ from $\ketbra{\sigma}{\rho}^{\sA'\sB}$.
This is achieved by a technique of the linear combination of block-encoded matrices~\cite{gilyen2019qsvt, Gilyn2019thesis} as follows: let $X_i$ be a tensor product of the Pauli-$X$ operators such that $X_i\ket{0} = \ket{i}$.
Then, a unitary
\begin{align}
\label{inteq:153}
    \til{U}^{\sA'\sB\sA''\sE} =  H^{\sA''}\Big(\sum_{i=1}^{d_\sA} \ketbra{i}{i}^{\sA''} \otimes X_i^{\sA'} U^{\sA'\sB\sE} X_i^{\sA'}\Big)H^{\sA''}, 
\end{align}
where $\sE = \sD_1\sD_2$, is a $(d_\sA^3d_\sB^2, \cO(\log{d_{\sA\sB}}), 2d_\sA^2 d_\sB\delta_1)$-block encoding of $\tr_{\sA'}[\ketbra{\sigma}{\rho}^{\sA'\sB}]$, because we have
\begin{align}
    &\big\|d_\sA^3 d_\sB^2 \bra{0}^{\sA'\sA''\sE} \til{U}^{\sA'\sB\sA''\sE} \ket{0}^{\sA'\sA''\sE} - \tr_{\sA'}\big[\ketbra{\sigma}{\rho}^{\sA'\sB}\big]\big\|_\infty \notag\\
    &= \big\|d_{\sA\sB}^2\sum_i\bra{i}^{\sA'}\bra{0}^\sE U^{\sA'\sB\sE}\ket{0}^\sE\ket{i}^{\sA'} - \tr_{\sA'}\big[\ketbra{\sigma}{\rho}^{\sA'\sB}\big]\big\|_\infty \\
    &= \big\|d_{\sA\sB}^2\tr_{\sA'}\big[\bra{0}^\sE U^{\sA'\sB\sE} \ket{0}^\sE\big]- \tr_{\sA'}\big[\ketbra{\sigma}{\rho}^{\sA'\sB}\big]\big\|_\infty \\
    &\leq d_\sA\big\|d_{\sA\sB}^2\bra{0}^\sE U^{\sA'\sB\sE} \ket{0}^\sE - \ketbra{\sigma}{\rho}^{\sA'\sB}\big\|_\infty \\
    \label{inteq:124}
    &\leq 2d_\sA^2 d_\sB\delta_1,
\end{align}
where the first inequality follows from the fact that, for any matrix $S^{\sA\sB}$, $\|\tr_\sA[S^{\sA\sB}]\|_\infty \leq d_\sA\|S^{\sA\sB}\|_\infty$~\cite{rastegin2012NormPartialTrace}, and the last inequality holds as $U^{\sA'\sB\sE}$ is a $(d_{\sA\sB}^2, \cO(\log{d_{\sA\sB}}), 2d_{\sA\sB}\delta_1)$-block-encoding of $\ketbra{\sigma}{\rho}^{\sA'\sB}$.
Since $\til{U}^{\sA'\sB\sA''\sE}$ in Eq.~\eqref{inteq:153} includes $d_\sA$ unitaries of the form $X_i^{\sA'} U^{\sA'\sB\sE} X_i^{\sA'}$ for $i = 1, 2, \ldots, d_\sA$, the number of queries up to this point is $l = d_\sA \til{\cO}(d_{\sA\sB}^6/\delta_1^2) = \til{\cO}(d_{\sA\sB}^7/\delta_1^2)$.

Let $\sF = \sA'\sA''\sE$.
Now, we see that the unitary $\til{U}^{\sB\sF}$ approximately takes the form:
\begin{align}
        \til{U}^{\sB\sF}
        \approx \hspace{1mm}
\begin{blockarray}{ccc}
 & \bra{0}^\sF &  & \vspace{1mm}\\
\begin{block}{c(cc)}
  \ket{0}^\sF \hspace*{1mm} & \f{1}{d_\sA^3 d_\sB^2}\tr_\sA\big[\ketbra{\sigma}{\rho}^{\sA\sB}\big] & \hspace{0mm} * \hspace{3mm}\\
   \hspace*{1mm} & * & \hspace{0mm} * \hspace{3mm}\\
\end{block}
\end{blockarray}\hspace{2.5mm}.
\end{align}
By applying the QSVT with the sign function, we construct a unitary $\til{W}^{\sB\sH}$ that block-encodes a degree-$u$ odd polynomial $P_\sgn^{(\rm SV)}(\bra{0}^\sF\til{U}^{\sB\sF}\ket{0}^\sF)$, using $\til{U}^{\sB\sF}$ for $u=\cO(\log(1/\delta_2)/\beta)$ times, where $\beta$ will be chosen appropriately later.
To evaluate the errors accumulated in this step, we use the following lemma~\cite{Gilyn2019thesis, gilyen2022improvedfidelity}.


\begin{lemma}[Robustness of the QSVT~{\cite[Lemma 2.4.4]{Gilyn2019thesis}}]
\label{lem:robstness of QSVT}
Suppose $P$ is an odd/even polynomial of degree-$l$ such that 
$|P(x)| \leq 1$ for $x\in[-1, 1]$, and suppose $A$ and $\tilde{A}$ are matrices with $\|A\|_\infty, \|\til{A}\|_\infty \leq 1$, such that $\|A-\til{A}\|_\infty + \Big\|\f{A+\til{A}}{2}\Big\|_\infty^2 \leq 1$. Then, we have 
\begin{align}
    \big\|P^{(\rm SV)}(A) - P^{(\rm SV)}(\til{A})\big\|_\infty &\leq l\sqrt{\f{2}{1-\big\|\f{A+\til{A}}{2}\big\|_\infty^2}}\big\|A-\til{A}\big\|_\infty,
\end{align}
where $P$ acts on the singular values of the input matrix.
\end{lemma}


Let $M^\sB = \tr_{\sA'}\big[\ketbra{\sigma}{\rho}^{\sA'\sB}\big]$. 
We observe that
\begin{align}
    \big\|P_\sgn^{(\rm SV)}\big(\bra{0}^\sF \til{U}^{\sB\sF}\ket{0}^\sF\big) - P_\sgn^{(\rm SV)}\big(M^\sB/(d_\sA^3d_\sB^2)\big)\big\|_\infty 
    &\leq 2u \big\|\bra{0}^\sF \til{U}^{\sB\sF}\ket{0}^\sF - M^\sB/(d_\sA^3 d_\sB^2)\big\|_\infty \\
    \label{inteq:134}
    &\leq \eta_1,
\end{align}
where $\eta_1 = 4u\delta_1/d_{\sA\sB}$.
In the second inequality, we used Eq.~\eqref{inteq:124}, and in the first inequality, we used Lemma~\ref{lem:robstness of QSVT} with
\begin{align}
    1 - \Big\|\f{1}{2}\Big(\bra{0}^\sF \til{U}^{\sB\sF}\ket{0}^\sF + M^\sB/(d_\sA^3 d_\sB^2)\Big)\Big\|_\infty^2 \geq 1/2,
\end{align}
which holds under the assumptions $\delta_1 \in (0, 1)$ and $d_\sA, d_\sB \geq 2$ that incurs no essential loss of generality.
Hence, we see that $\til{W}^{\sB\sH}$ is a $(1, \cO(\log{d_{\sA \sB}}), \eta_1)$-block-encoding unitary of $P_\sgn^{(\rm SV)}\big(M^\sB/(d_\sA^3d_\sB^2)\big)$.



We finally evaluate the approximation error using the fidelity difference.
Let
\begin{align}
    \til{\rF} 
    &= \rF\big(\til{W}^{\sB\sH}\ket{\rho}^{\sA\sB}\ket{0}^\sH, \ket{\sigma}^{\sA\sB}\ket{0}^\sH\big) \\
    &= \big|\bra{\sigma}^{\sA\sB}P_\sgn^{(\rm SV)}\big(\bra{0}^\sF \til{U}^{\sB\sF}\ket{0}^\sF\big)\ket{\rho}^{\sA\sB}\big|,
\end{align}
$\rF' = \big|\bra{\sigma}^{\sA\sB}P_\sgn^{(\rm SV)}\big(M^\sB/(d_\sA^3d_\sB^2)\big)\ket{\rho}^{\sA\sB}\big|$, and $\rF = \rF(\rho^\sA, \sigma^\sA) = \big|\bra{\sigma}^{\sA\sB} V^{\sB}\ket{\rho}^{\sA\sB}\big|$,
where $V^\sB$ is the exact Uhlmann partial isometry.
Note that $\sgn^{(\rm SV)}\big(M^\sB/(d_\sA^3d_\sB^2)\big)$ corresponds to $V^\sB$.
Then, we compute the differences between them as
\begin{align}
    \Big|\sqrt{\til{\rF}} - \sqrt{\rF'}\Big| 
    &\leq \Big|\tr\Big[\ketbra{\rho}{\sigma}^{\sA\sB}\Big(P_\sgn^{(\rm SV)}\big(\bra{0}^\sF \til{U}^{\sB\sF}\ket{0}^\sF\big) - P_\sgn^{(\rm SV)}\big(M^\sB/(d_\sA^3d_\sB^2)\big)\Big)\Big]\Big| \\
    &\leq \|\ketbra{\rho}{\sigma}^{\sA\sB}\|_1 \big\|P_\sgn^{(\rm SV)}\big(\bra{0}^\sF \til{U}^{\sB\sF}\ket{0}^\sF\big) - P_\sgn^{(\rm SV)}\big(M^\sB/(d_\sA^3d_\sB^2)\big)\big\|_\infty \\
    \label{inteq:125}
    &\leq \eta_1,
\end{align}
where in the second inequality we used
$|\tr[AB]| \leq \|AB\|_1 \leq \|A\|_1\|B\|_\infty$, which follows from the H\"{o}lder's inequality in Eq.~\eqref{eq:Holder ineq}, and in the last line we used Eq.~\eqref{inteq:134} and $\|\ketbra{\rho}{\sigma}^{\sA\sB}\|_1 = 1$. 

Also, from a calculation similar to that in Sec.~\ref{sec:uhlfidquery}, we see that
\begin{align}
    \big|\sqrt{\rF'} - \sqrt{\rF}\big| 
    &\leq \Big|\tr\Big[\Big(\sgn^{(\rm SV)}\big(M^\sB/(d_\sA^3 d_\sB^2)\big) - P_\sgn^{(\rm SV)}\big(M^\sB/(d_\sA^3 d_\sB^2)\big)\Big)(M^{\sB})^\dag\Big]\Big| \\
    &\leq \sum_k \big|1 - P_\sgn\big(s_k/(d_\sA^3 d_\sB^2)\big)\big|s_k \\
    &= \sum_{k\in I_\beta} \big|1 - P_\sgn\big(s_k/(d_\sA^3 d_\sB^2)\big)\big|s_k + \sum_{k\in \bar{I}_\beta} \big|1 - P_\sgn\big(s_k/(d_\sA^3 d_\sB^2)\big)\big|s_k \\
    &\leq \delta_2 + 2\beta r,
\end{align}
where $\{s_k\}_k$ and $r$ are the singular values and the rank of $M^\sB$, respectively, and $I_\beta = \{k \in \bN; s_k/(d_\sA^3 d_\sB^2) \geq \beta\}$ and $\bar{I}_\beta = \{k \in \bN; s_k/(d_\sA^3 d_\sB^2) < \beta\}$.
Here, we used $\big|\sgn(x) - P_\sgn(x)\big| \leq \delta_2$ for $x \in [\beta, 1]$ when $u=\cO\big(\log(1/\delta_2)/\beta\big)$.
Thus, setting $\beta$ as $1/\beta = \min\{d_\sA^3 d_\sB^2/s_{\rm min}, 2r/\delta_2\}$ ensures at least 
\begin{align}
\label{inteq:126}
    \big|\sqrt{\rF} - \sqrt{\rF'}\big| \leq 2\delta_1.
\end{align}

From Eqs.~\eqref{inteq:125} and~\eqref{inteq:126}, we obtain that $\big|\til{\rF} - \rF\big| \leq 2\eta_1 + 4\delta_2$, where we used the inequality that $|x-y|\leq 2 |\sqrt{x} - \sqrt{y}|$ for $0 \leq x, y \leq 1$.
Hence, when we define $\cT^\sB$ as $\cT^\sB = \tr_\sH \circ \cW^{\sB\sH} \circ \cP_{\ket{0}}^{\bC\rarr\sH}$, we have that
\begin{align}
    \rF\big(\cT^\sB(\ketbra{\rho}{\rho}^{\sA\sB}), \ket{\sigma}^{\sA\sB}\big) \geq \rF(\rho^\sA, \sigma^\sB) -2 \eta_1 - 4\delta_2.
\end{align}
By rescaling the parameters as $\eta_1 = 4u\delta_1/d_{\sA\sB} = \delta/4$, i.e., $\delta_1 = d_{\sA\sB}\delta/(16u)$, and $\delta_2 = \delta/8$, we finally obtain
\begin{align}
    \rF\big(\cT^\sB(\ketbra{\rho}{\rho}^{\sA\sB}) - \ket{\sigma}^{\sA\sB}\big) \geq \rF(\rho^\sA, \sigma^\sB) - \delta.
\end{align}
Consequently, the total number of queries in the entire algorithm is given by
\begin{align}
    \gamma 
    &= lu \\
    &= \til{\cO}\big(d_{\sA\sB}^7/\delta_1^2\big)\cO\big(\log{(1/\delta_2)}/\beta\big) \\
    &= \til{\cO}\Big(\f{d_{\sA\sB}^5}{\delta^2}\min\Big\{\f{d_\sA^9 d_\sB^6}{s_{\rm min}^3}, \f{r^3}{\delta^3}\Big\}\Big).
\end{align}
The most foundational operation throughout this algorithm is the Pauli measurement, which, up to logarithmic factors, has to be performed as many times as the total number of queries. 
Thus, the quantum circuit of this algorithm consists of $\tilde{\cO}(\gamma)$ one- and two-qubit gates.
Due to its sequential structure, $\cO(\log d_{\sA\sB})$ qubits suffice at any one time.






%===============================================================================================================================================================================================================







\section{Derivation of an inequality between the minimum non-zero singular values}
\label{sec:appendix product of singvalue}


We here derive the following proposition. We use $\supp[A]$ to denote the support of $A$, i.e., the orthogonal complement of the kernel of $A$, and $\im[A]$ to denote the image of $A$.

\begin{proposition}[Inequality about the minimum non-zero singular value of product of matrices]
\label{prop:min singular inequality}
For any matrices $A$ and $B$ such that $\im[B] \subset \supp[A]$ or $\im[A] \subset \supp[B]$, it holds that $s_{\rm min}(AB) \geq s_{\rm min}(A) s_{\rm min}(B)$, where $s_{\rm min}(\cdot)$ denotes the minimum non-zero singular value of the input matrix.
\end{proposition}

Note that if a matrix $A$ is Hermitian, then $\im[A] = \supp[A]$ holds.




\begin{proof}[Proof of \Cref{prop:min singular inequality}]


Since $s_{\rm min}(AB) = s_{\rm min}(BA)$, without loss of generality, we focus on the case where $\im[B] \subset \supp[A]$.
From the min-max principle~\cite{Gohberg2003basicclass, Simon2005traceideals, bhatia2013matrix}, the minimum non-zero singular value of $A$ is given by $s_{\rm min}(A) = \min_{x \in \supp[A]; x \neq 0}\f{\|Ax\|}{\|x\|}$. Hence, $s_{\rm min}(AB)$ can be calculated as follows:
\begin{align}
    s_{\rm min}(AB) 
    &= \min_{x \in \supp[AB]; x \neq 0}\f{\|ABx\|}{\|x\|} \\
    &= \min_{x \in \supp[AB]; x \neq 0}\f{\|ABx\|}{\|Bx\|}\f{\|Bx\|}{\|x\|} \\
    &\geq \min_{x \in \supp[AB]; x \neq 0}\f{\|ABx\|}{\|Bx\|}\min_{x \in \supp[AB]; x \neq 0}\f{\|Bx\|}{\|x\|} \\
    &\geq \min_{y \in \supp[A]; y \neq 0}\f{\|Ay\|}{\|y\|}\min_{x \in \supp[B]; x \neq 0}\f{\|Bx\|}{\|x\|} \\
    &= s_{\rm min}(A) s_{\rm min}(B),
\end{align}
where, in the second inequality, we used the fact that $\supp[AB] \subset \supp[B]$ and the assumption $\im[B] \subset \supp[A]$.  


\end{proof}




%=================================================================================================================================================================================================================================================



\section{Optimal local operation and the Uhlmann transformation}
\label{sec:opt local and Uhlmann}


We show that the Uhlmann's theorem characterizes an optimal local transformation between two states, even when the involved states are generally mixed states.

\begin{proposition}[Optimal local operation via the Uhlmann transformation]
\label{prop:opt local mixed via Uhl}
Let $\ket{\rho}^{\sA\sB\sE}$ and $\ket{\sigma}^{\sA\sC\sE}$ be purified states of $\rho^{\sA\sB}$ and $\sigma^{\sA\sC}$, respectively. Then, any quantum channel $\cT^{\sB\rarr\sC}$ satisfies that
    \begin{align}
    \label{eq:opt ineq local op}
        \rF\big(\cT^{\sB\rarr\sC}(\rho^{\sA\sB}), \sigma^{\sA\sC}\big) \leq \max_{\cL^\sE} \rF\big(\rho^{\sA\sE}, \cL^\sE(\sigma^{\sA\sE})\big),
    \end{align}
where maximization is taken over all quantum channels acting on the system $\sE$.

Let $\til{\cL}^\sE$ be a quantum channel that attains the maximization in Eq.~\eqref{eq:opt ineq local op}, and $V_{\til{\cL}}^{\sE\rarr\sE\sF}$ be its Stinespring isometry.
A quantum channel $\til{\cT}^{\sB\rarr\sC}$ that achieves equality in Eq.~\eqref{eq:opt ineq local op} is given by
\begin{align}
\til{\cT}^{\sB\rarr\sC}(\cdot) = \tr_\sF\big[U^\sG(\cdot \otimes \ketbra{0}{0}^\sD)(U^\sG)^\dag\big],
\end{align}
where $U^\sG$ is an Uhlmann unitary acting on $\sG = \sB\sD = \sC\sF$, which satisfies 
\begin{align}
    \rF\big(U^\sG\ket{\rho}^{\sA\sB\sE}\ket{0}^\sD, V_{\til{\cL}}^{\sE\rarr\sE\sF}\ket{\sigma}^{\sA\sC\sE}\big) = \rF\big(\rho^{\sA\sE}, \til{\cL}^\sE(\sigma^{\sA\sE})\big).
\end{align}
\end{proposition}




Note that, since the optimal transformation $\til{\cT}^{\sB\rarr\sC}$ is constructed from the Uhlmann unitary $U^\sG$ connecting $\ket{\rho}^{\sA\sB\sE}\ket{0}^\sD$ and $V_{\til{\cL}}^{\sE\rarr\sE\sF}\ket{\sigma}^{\sA\sC\sE}$, implementing it would generally require knowledge of the purified states $\ket{\rho}^{\sA\sB\sE}$, $\ket{\sigma}^{\sA\sB\sE}$, and also the quantum channel $\til{\cL}^\sE$ that attains the maximization of Eq.~\eqref{eq:opt ineq local op}.

In particular, when we apply this proposition to the decoupling argument discussed in Sec.~\ref{sec:decoupling and uhlmann}, the state $\sigma^{\sA\sE}$ can be taken to be a product state $\rho^\sA \otimes \tau^\sE$. In this case, the right-hand side of Eq.~\eqref{eq:opt ineq local op} reduces to
$\max_{\tau^\sE}\rF(\rho^{\sA\sE}, \rho^\sA\otimes\tau^\sE)$, where the maximization is taken over all states on $\sE$.
By identifying a state $\til{\tau}^\sE$ that achieves this maximization, we can construct the optimal local operation via the Uhlmann transformation.

We now prove the proposition.


\begin{proof}[Proof of \Cref{prop:opt local mixed via Uhl}]
    


From the Uhlmann's theorem, for any quantum channel $\cT^{\sB\rarr\sC}$, there is a unitary $\hat{U}^{\sE\sF}$ such that
\begin{align}
\label{inteq:108}
    \rF\big(\cT^{\sB\rarr\sC}(\rho^{\sA\sB}), \sigma^{\sA\sC}\big) = \rF\big(V_\cT^{\sB\rarr\sC\sF}\ket{\rho}^{\sA\sB\sE}, \hat{U}^{\sE\sF}\ket{\sigma}^{\sA\sC\sE}\ket{0}^\sF\big),
\end{align}
where $V_\cT^{\sB\rarr\sC\sF}$ is a Stinespring isometry of $\cT^{\sB\rarr\sC}$.
Since the fidelity is monotonic under the partial trace, the right-hand side on Eq.~\eqref{inteq:108} is bounded as
\begin{align}
    \rF\big(V_\cT^{\sB\rarr\sC\sF}\ket{\rho}^{\sA\sB\sE}, \hat{U}^{\sE\sF}\ket{\sigma}^{\sA\sC\sE}\ket{0}^\sF\big)
    &\leq \rF\big(\rho^{\sA\sE}, \tr_\sF\big[\hat{U}^{\sE\sF}(\sigma^{\sA\sE}\otimes\ketbra{0}{0}^\sF)(\hat{U}^{\sE\sF})^\dag\big]\big) \\
    &\leq \max_{\cL^{\sE}}\rF\big(\rho^{\sA\sE}, \cL^\sE(\sigma^{\sA\sE})\big), 
\end{align}
where we take the maximization over all quantum channels acting on system $\sE$, considering that $\tr_\sF\big[\hat{U}^{\sE\sF}(\cdot\otimes\ketbra{0}{0}^\sF)(\hat{U}^{\sE\sF})^\dag\big]$ gives a quantum channel on $\sE$.
Thus, any quantum channel $\cT^{\sB\rarr\sC}$ satisfies
\begin{align}
\label{inteq:109}
    \rF\big(\cT^{\sB\rarr\sC}(\rho^{\sA\sB}), \sigma^{\sA\sC}\big)
    \leq \max_{\cL^{\sE}}\rF\big(\rho^{\sA\sE}, \cL^\sE(\sigma^{\sA\sE})\big).
\end{align}

Regarding the equality condition, let $\til{\cL}^\sE$ be a quantum channel that attains the maximization, and let $V_{\til{\cL}}^{\sE \to \sE\sF}$ be its Stinespring isometry.
Then, from the Uhlmann's theorem, there exists a unitary $U^\sG$ on $\sG = \sB\sD = \sC\sF$ satisfying that
\begin{align}
    \rF\big(\rho^{\sA\sE}, \til{\cL}^\sE(\sigma^{\sA\sE})\big)
    &= \rF\big(U^\sG\ket{\rho}^{\sA\sB\sE}\ket{0}^\sD, V_{\til{\cL}}^{\sE\rarr\sE\sF}\ket{\sigma}^{\sA\sC\sE}\big) \\
    \label{inteq:110}
    &\leq \rF\big(\til{\cT}^{\sB\rarr\sC}(\rho^{\sA\sB}), \sigma^{\sA\sC}\big),
\end{align}
where $\til{\cT}^{\sB\rarr\sC}(\cdot) = \tr_\sF\big[U^\sG(\cdot\otimes\ketbra{0}{0}^\sD)(U^\sG)^\dag\big]$, and in the second inequality we used the monotonic property of the fidelity under the partial trace again.
From Eqs.~\eqref{inteq:109} and~\eqref{inteq:110}, we see that $\til{\cT}^{\sB\rarr\sC}$ gives the optimal transformation from $\rho^{\sA\sB}$ to $\sigma^{\sA\sC}$ through a local operation and achieves that $\rF\big(\til{\cT}^{\sB\rarr\sC}(\rho^{\sA\sB}), \sigma^{\sA\sC}\big) = \max_{\cL^{\sE}}\rF\big(\rho^{\sA\sE}, \cL^\sE(\sigma^{\sA\sE})\big)$.

\end{proof}




%===========================================================================================================================================


